\documentclass[11pt]{article}
\usepackage{fontspec}
\usepackage[french]{babel}
\usepackage{amsmath}
\usepackage{amssymb}
\usepackage{amsthm}
\usepackage{mathrsfs}
\usepackage{tikz-cd}
\usepackage{enumitem}

\newcommand{\printedpage}[1]{\marginpar{\footnotesize p.~#1}}
\newcommand{\msp}[1]{\marginpar{\footnotesize [#1]}}
\newcommand{\Filt}{\mathrm{Filt}}
\newcommand{\card}{\operatorname{card}}
\newcommand{\Ob}{\operatorname{Ob}}
\newcommand{\Fl}{\operatorname{Fl}}
\newcommand{\Ind}{\mathrm{Ind}}
\newcommand{\cardred}{\text{card réd}\,}

\begin{document}

\printedpage{54}
% head: A. Grothendieck : Dérivateurs, Ch. XVIII, Édition des Universités de Montpellier II et Paris VII

\textsc{Démonstration} de 4.7.1. Les implications a $\Longrightarrow$ a$'$ $\Longrightarrow$ a$''$ sont tautologiques, ainsi que b $\Longrightarrow$ b$'$. L’implication b$'$ $\Longrightarrow$ c$'$ provient de la stabilité de $C^{\wedge}_{\pi}$ par $\varinjlim$ de cardinal $\leqslant \pi$, et du fait que les $\varinjlim$ du type envisagé (en particulier, grandes devant $\pi_0$) \textit{dans} $\mathcal{M}$ sont les mêmes que dans $C^{\wedge}$, i.e. se calculent “argument par argument” (cf. 4.4.15, p. 99). Ceci, plus le fait que l’inclusion $\mathcal{M} \hookrightarrow C^{\wedge}$ est pleinement fidèle, implique aussi c$'$ $\Longrightarrow$ c.

\msp{page 164}
L’implication a$''$ $\Longrightarrow$ c$'$ se voit comme dans 4.5.6 (p. 119) en appliquant le lemme 4.5.8 1° et 2°, au cas de $C^{\wedge} = \underline{\mathrm{Hom}}(C^{\circ}, \mathcal{E})$, où $\mathcal{E} =$ (Ens). L’implication conditionnelle inverse c$'$ $\Longrightarrow$ a$''$, pour $\pi \geqslant \pi_1$, se voit comme dans loc. cit., en appliquant ce lemme 4.5.8. Il reste donc à établir encore les trois implications
\[
\left\{
\begin{array}{ccll}
\text{a} & \Longrightarrow & \text{b} & \\[4pt]
\text{a}'' & \Longrightarrow & \text{a} & \text{si } \pi \geqslant \pi_1 \\[4pt]
\text{c} & \Longrightarrow & \text{b} & \text{si } \pi = \pi^{\pi_0}.
\end{array}
\right. \tag*{$(*)$}
\]
Mais a $\Longrightarrow$ b est tautologique, en prenant $I = C_0/F$, puisqu’on sait que $C_0/F$ est grande devant $\pi_0$ et que $F = \varinjlim\limits_{i \in I} x_i$, pratiquement par définition de $\mathcal{M} = \underbrace{\Ind_{\pi}(C)}_{\simeq \Ind_{\pi} C_0} \subset C^{\wedge}$. Donc il reste à établir les deux implications conditionnelles de $(*)$. L’implication a$''$ $\Longrightarrow$ a pour $\pi \geqslant \pi_1$ s’établit comme dans 4.5.2 (cf. page 109), c’est essentiellement sorital. Reste à établir
\[
\text{c} \;\Longrightarrow\; \text{b} \qquad \text{si} \;\; \pi \;=\; \pi^{\pi_0} \;. \tag*{$(*)$}
\]
Nous aurons besoin du résultat auxiliaire :

\msp{page 165}
\textbf{Lemme 4.7.3.} \textit{Soit $\pi \geqslant \pi_0$ tel que $\pi^{\pi_0} = \pi$ (donc en fait $\pi > \pi_0$). Pour tout objet $F$ de $\mathcal{M}$, on a}
\[
F \;=\; \varinjlim_{i \in I} x_i \qquad \text{(\footnotemark[67])}\;. \tag*{(4.7.3.1)}
\]
\footnotetext[67]{les $x_i$ dans $\Filt^{\pi}$, $I$ ensemble ordonné grand devant $\pi$.}
\textit{Si de plus $\pi' > \pi$ et $F \in \Ob \Filt^{\pi'}$, alors on peut prendre $I$ telle que}
\[
\card I \;\leqslant\; \pi'^{\pi} \;. \tag*{(4.7.3.2)}
\]

En effet, on a
\[
F \;=\; \varinjlim_{j \in J} x_j \qquad \text{avec } J = C_0/F \qquad \text{(\footnotemark[68])}\;,
\]
\footnotetext[68]{où $J$ peut aussi se remplacer par une catégorie cofinale dans [?] $C_0/F$ grande devant $\pi_0$.}
avec $J$ grande devant $\pi_0$. Soit alors
\[
I \;\;=\;\; \parbox[t]{0.62\textwidth}{ensemble des sous-catégories $J'$ de $J$ \textit{grandes devant} $\pi_0$, et telles que $\card J' \leqslant \pi$.} \tag*{(4.7.3.3)}
\]

\printedpage{55}
% head: A. Grothendieck : Dérivateurs, Ch. XVIII, Édition des Universités de Montpellier II et Paris VII

Je dis que $I$ (ordonné par inclusion) \textit{est grand devant} $\pi$ (\footnote[69]{cf. démonstration 4.7.5, page 168-171.}). Cela résulte aussitôt du lemme 4.5.5, p. 113, (énoncé d’abord, de façon erronée, en supposant seulement $\pi > \pi_0$ et non $\pi^{\pi_0} = \pi$). Posons alors, pour tout $J' \in I$,
\[
x_{J'} \;=\; \varinjlim_{j \in J'} x_j \;,
\]
on a donc
\[
x_{J'} \;\in\; \Ob \Filt^{\pi}
\]
par définition de $I$ et de $\Filt^{\pi}$. D’autre part,
\[
J' \;\longmapsto\; x_{J'} \;:\; I \;\longrightarrow\; \Filt^{\pi}
\]
\msp{page 166}
est un foncteur, donc
\[
\varinjlim_{J' \in I} x_{J'} \;\simeq\; \varinjlim_{j \in J} x_j \;\simeq\; x \;, % [sic] x -> F, the colimit is the object F of (4.7.3.1)
\]
d’où la première assertion de 4.7.3. Reste la majoration de cardinaux pour $F \in \Ob \Filt^{\pi}$. Dans ce cas, quitte à remplacer $C_0/F$ par une catégorie cofinale grande devant $\pi_0$, $\card J \leqslant \pi'$. On voit alors tout de suite que $\card I \leqslant \pi'^{\pi_0}$, ce qui est la majoration annoncée.

Nous pouvons maintenant prouver l’implication $(*)$, page 164, i.e.
\[
\mathcal{M}_{\pi} \;\subset\; \Filt^{\pi} \qquad (\text{pour } \pi = \pi^{\pi_0}).
\]
Soit donc $F$ dans $\mathcal{M}_{\pi}$, i.e. $\pi$-accessible, écrivons $F$ sous la forme (4.7.3.2), avec $I$ grande devant $\pi$. Par définition de la $\pi$-accessibilité, appliqué à $F \xrightarrow{\mathrm{id}} F \simeq \varinjlim F_i$, on voit que $F$ est un facteur direct d’un $F_i$, donc d’un objet de $\Filt^{\pi}$. Faute de savoir que $\Filt^{\pi}$ est stable par facteurs directs, on ne peut en conclure que $F$ est dans $\Filt^{\pi}$ (aussi je ne peux prouver c $\Longrightarrow$ b), mais seulement que $F$ est dans $\Filt'^{\pi}$, qui, lui, est stable par $\varinjlim$ grandes devant $\pi_0$, de cardinal $\pi$ donné $> \pi_0$, donc aussi \msp{page 167} stable par facteurs directs. Cela achève la démonstration de 4.7.1, et en même temps on obtient :

% marginal vertical bar beside Corollaire 4.7.4 and its statement
\textbf{Corollaire 4.7.4.} \textit{Si $\pi \geqslant \pi_0$ est tel que $\pi^{\pi_0} = \pi$, alors on a}
\[
\Filt'^{\pi} \;=\; (\Filt^{\pi})^{\natural} \;, \tag*{(4.7.4.1)}
\]
\textit{i.e.} $\Filt'^{\pi}$ \textit{est formée des objets $F$ de $\mathcal{M}$ qui sont facteurs directs d’un objet de} $\Filt^{\pi}$.

En effet, comme $\Filt^{\pi} \subset \Filt'^{\pi}$, et que $(\Filt^{\pi})^{\natural} \subset \Filt'^{\pi}$, et que $\Filt'^{\pi}$ est stable par facteurs directs, on a $(\Filt^{\pi})^{\natural} \subset \Filt'^{\pi}$ pour tout $\pi \geqslant \pi_0$. Mais inversement, comme $\Filt'^{\pi} \subset \mathcal{M}_{\pi}$, l’argument précédent montre que pour $\pi = \pi^{\pi_0}$, on a $\Filt'^{\pi} \subset \mathcal{M}_{\pi} \subset (\Filt^{\pi})^{\natural}$, ce qui prouve 4.7.3.

\textbf{Remarque 4.7.4.1} \texttt{[deux fois 4.7.4.1]}. Bien sûr, par 4.7.2, 1°, dès que $\pi \geqslant \pi_1$, on a même $\Filt^{\pi} = \Filt'^{\pi}$ (sans avoir à supposer $\pi = \pi^{\pi_0}$), donc \textit{si} $\pi > \pi_0$ (ce qui sera la cas si $\pi_1 > \pi_0$, i.e. $\pi_1 \neq \pi_0$). $\Filt'^{\pi}$ étant stable par facteurs directs, $\Filt^{\pi}$ l’est aussi, (4.7.4.1) est satisfait trivialement. Mais si $\pi = \pi_1 = \pi_0$, on a $\Filt^{\pi} = \Filt^{\pi_0} = \Filt'^{\pi_0} =$ sous-catégorie des foncteurs \textit{représentables} dans $C$ (équivalent à $C$), et celle-ci n’est égale à $C^{\natural}$ dans $\mathcal{M}$ que si $C$ est karoubienne, cf. \ldots{} ci-dessous. % [sic] la cas -> le cas, gender of "cas"

\printedpage{56}
% head: A. Grothendieck : Dérivateurs, Ch. XVIII, Édition des Universités de Montpellier II et Paris VII

\msp{page 168}
\textbf{4.7.5.} Pour que la démonstration de 4.7.1 soit complète, il me faudrait encore donner une démonstration de 4.5.5, p. 113, qui est le lemme technique clef pour établir le

\textbf{Lemme 4.7.5.1} (utilisé dans la démonstration de 4.7.3). \textit{Soient $\pi_0$, $\pi$ deux cardinaux infinis, avec $\pi \geqslant \pi_0$ et $\pi^{\pi_0} = \pi$ (donc $\pi > \pi_0$). Soit $J$ une catégorie grande devant $\pi_0$. Alors l’ensemble $I$, ordonné par inclusion, des sous-catégories $J'$ de $J$ qui sont grandes devant $\pi_0$, et telles que} $\card J' \leqslant \pi$, \textit{est grand devant $\pi$. (De plus, si} $\card J \leqslant \pi'$, \textit{on a} $\card I \leqslant \pi'^{\pi}$\textit{.)}

Il faut donc montrer le

\textbf{Corollaire 4.7.5.2} (\footnote[70]{Comparer avec la démonstration de 4.6.1.2.}) (c’est le lemme 4.5.5 non encore démontré). \textit{Sous les conditions de 4.7.5, soit $K \subset \mathrm{Fl} J$ tel que} $\card K \leqslant \pi$. \textit{Alors il existe $J' \in I$ tel que $K \subset J'$.}

On sait, $\pi$ étant infini, que la catégorie $K'$ engendrée par $K$ est aussi de cardinal $\leqslant \pi$, donc on peut supposer que $K$ “est” une sous-catégorie de $J$. On va construire une $J' \in I$ contenant $K$, en construisant par récurrence transfinie une suite croissante $K_{\alpha}$ de sous-catégories de $J$, avec $\card K_{\alpha} \leqslant \pi$, indexée par les ordinaux $\alpha$ tels que \fbox{$\card \alpha \leqslant \pi$} :

\msp{page 169}
\begin{enumerate}
\item[1°)] $K_0 = K$.
\item[2°)] $K_{\alpha} = \bigcup_{\beta < \alpha} K_{\beta}$ ($= \varinjlim K_{\beta}$) si $\alpha$ est un ordinal limite (donc $\card K_{\alpha} \leqslant \pi$, puisque les $\card K_{\beta} \leqslant \pi$).
\item[3°)] $K_{\alpha+1}$ déduit de $K_{\alpha}$ de la façon suivante. Si $K_{\alpha}$ est grande devant $\pi_0$, on prend $K_{\alpha+1} = K_{\alpha}$ (donc, compte tenu de 2°, on aura $K_{\alpha'} = K_{\alpha}$ pour $\alpha' \geqslant \alpha$ -- la construction transfinie s’arrête \ldots). Sinon, $K_{\alpha+1}$ s’obtient ainsi :
\begin{enumerate}
\item[a)] pour toute partie $A$ de cardinal $\leqslant \pi_0$ de $\Ob K_{\alpha}$, on choisit un majorant $x_A$ dans $J$ (on a envie de le prendre dans $K_{\alpha}$, s’il y a un majorant dans $K_{\alpha}$)\,;
\item[b)] pour tout couple $x$, $y$ d’objets de $K_{\alpha}$, et toute partie de cardinal $\leqslant \pi_0$ de $\mathrm{Hom}_{K_{\alpha}}(x,y)$, on choisit une flèche $u_B : y \longrightarrow z_B$ dans $J$, qui égalise les $v_{\beta} \in B$ (et on prend cette flèche dans $K_{\alpha}$, quand c’est possible)\,;
\item[c)] on prend pour $K_{\alpha+1}$ la plus petite sous-catégorie contenant les $\mathrm{id}_{x_A}$ et les $u_B$.
\end{enumerate}
\end{enumerate}

Il faut vérifier que
\[
\card K_{\alpha+1} \;\leqslant\; \pi \;.
\]
Le nombre des choix possibles pour $A \in \mathfrak{P}_{\leqslant \pi_0}(\Ob J)$ \msp{page 170} est $\leqslant \pi^{\pi_0} = \pi$, le nombre des choix possibles pour $B \in \mathfrak{P}_{\pi_0}(\mathrm{Fl}\, J)$ est aussi $\leqslant \pi^{\pi_0} = \pi$, donc le nombre des flèches de la forme $\mathrm{id}_{x_A}$ ou $u_B$ est $\leqslant \pi + \pi = \pi$. Mais alors la sous-catégorie $K_{\alpha+1}$ de $J$ engendrée par cet ensemble de flèches est de cardinal $\leqslant 1$. Soit alors % [sic] \leqslant 1 -> \leqslant \pi, the statement to be verified is card K_{\alpha+1} \leqslant \pi
\printedpage{57}
% head: A. Grothendieck : Dérivateurs, Ch. XVIII, Édition des Universités de Montpellier II et Paris VII
$\alpha_0$ le plus petit ordinal tel que
\[
\card \alpha_0 \;>\; \pi_0 \;,
\]
on a donc
\[
\card \alpha_0 \;=\; \text{cardinal successeur de } \pi_0 \leqslant \pi \;,
\]
puisque $\pi^{\pi_0} = \pi$, d’où $\pi > \pi_0$ (car $\pi^{\pi_0} = \pi_0^{\pi_0} \geqslant 2^{\pi_0} > \pi_0$). Ainsi, $\alpha_0$ fait partie des ordinaux $\alpha$ admissibles (i.e. tels que $\card \alpha \leqslant \pi$). Je dis que $K_{\alpha_0}$ \textit{est grande devant} $\pi_0$ (ce qui prouvera 4.7.5.2, en prenant $J' = K_{\alpha_0}$). Cela résulte immédiatement du fait que
\[
K_{\alpha_0} \;=\; \bigcup_{\alpha < \alpha_0} K_{\alpha} \qquad (\text{réunion \textit{croissante}}),
\]
% double marginal vertical bar beside the following two lines
et que l’ensemble des $\alpha < \alpha_0$ est \textit{grande devant} $\pi_0$ (\footnote[71]{expliciter ceci en un

\textbf{Lemme.} \textit{Soit $\pi_0$ un cardinal, $\alpha_0 = \alpha_{\pi'_0}$ ($\pi'_0$ le successeur de $\pi_0$) le plus petit ordinal tel que} $\card \alpha_0 > \pi_0$, \textit{alors $I_{\alpha_0} = \{\alpha \mid \alpha < \alpha_0\}$ est grand devant $\pi_0$.}}). Donc tout ensemble d’objets ou de flèches de $K_{\alpha_0}$, de cardinal $\leqslant \pi_0$, est contenu dans un $K_{\alpha}$ avec $\alpha < \alpha_0$.

\msp{page 171}
Cela implique, s’il s’agit d’un ensemble d’objets, qu’il y a un majorant commun dans $K_{\alpha+1}$, a fortiori dans $K_{\alpha_0}$. S’il s’agit d’une partie d’un $\mathrm{Hom}_{K_{\alpha}}(x,y)$, qu’il y a un égalisateur $y \longrightarrow z$ dans $K_{\alpha+1}$, a fortiori dans $K_{\alpha_0}$. Donc le critère $\mathrm{PF}_{\pi}$ 1,2 de 4.4.4 (page 89) est satisfait par $K_{\alpha_0}$, ce qui achève la démonstration de 4.7.5.2.

\textbf{Corollaire 4.7.6} (de 4.7.1). \textit{Pour tout cardinal $\pi \geqslant \pi_0$, on a}
\[
\underbrace{\Filt'^{\pi}}_{\overset{\textit{si } \pi \geqslant \pi_1}{=}\ \Filt^{\pi}} \;\subset\; \mathcal{M}_{\pi} \;\subset\; \underbrace{\Filt^{\pi^{\pi_0}}}_{\overset{\textit{si } \pi^{\pi_0} \geqslant \pi_1}{=}\ \Filt^{\pi^{\pi_0}}} \qquad (= (\Filt^{\pi^{\pi_0}})^{\natural} \textit{ par cor. 4.7.4})\,, \tag*{(4.7.6.1)} % ?? uncertain: tiny condition under the second brace -- "si \pi^{\pi_0} \geqslant \pi_1" or "si \pi^{\pi}_0 \geqslant \pi_1"
\]
\textit{donc si $\pi \geqslant \pi_1$, on a aussi}
\[
\Filt^{\pi} \;\subset\; \mathcal{M}_{\pi} \;\subset\; \Filt^{\pi^{\pi_0}} \qquad \textit{(si $\pi^{\pi_0} \geqslant \pi_1$).} \tag*{(4.7.6.2)}
\]

En effet, $\Filt'^{\pi} \subset \mathcal{M}_{\pi}$ n’est autre que l’implication b$'$ $\Longrightarrow$ c de 4.7.1, d’autre part, on a $\pi \leqslant \pi^{\pi_0}$, d’où
\[
\mathcal{M}_{\pi} \;\subset\; \mathcal{M}_{\pi^{\pi_0}} \;=\; \Filt'^{\pi^{\pi_0}} \;,
\]
où la dernière égalité provient de $\mathcal{M}_{\pi'} = \Filt'^{\pi'}$ si $\pi'^{\pi_0} = \pi'$, i.e. de (4.7.2.3) (page 163), en faisant $\pi' = \pi^{\pi_0}$, compte tenu que $\pi'^{\pi_0} = (\pi^{\pi_0})^{\pi_0} = \pi^{\pi_0 \cdot \pi_0} = \pi^{\pi_0} = \pi$.

\msp{page 172}
\textbf{Proposition 4.7.7.} \textit{Considérons $C$, $C_0$, $\mathcal{M} = \Ind_{\pi_0}(C)$, $\pi_1 = \sup(\pi_0, \card \mathrm{Fl}\, C_0)$ comme précédemment. Pour tout cardinal $\pi \geqslant \pi_0$, soit $\mathcal{F}^{\pi}$ la sous-catégorie strictement pleine} $\Filt'^{\pi}$ \textit{(resp. $(C^{\wedge})_{\pi} \cap \mathcal{M}$, resp. $\mathcal{M}_{\pi}$) dans $\mathcal{M}$. Alors la filtration croissante $(\mathcal{F}^{\pi})_{\pi \geqslant \pi_0}$ de $\mathcal{M}$ par les sous-catégories $\mathcal{F}^{\pi}$ satisfait les conditions suivantes.}
\begin{enumerate}
\item[\textit{a)}] \textit{Pour tout $\pi \geqslant \pi_0$, $\mathcal{F}^{\pi}$ est équivalente à une petite catégorie.}
\item[\textit{b)}] \textit{$\mathcal{M}$ satisfait à $\mathrm{L}_{\pi_0}$ (i.e. est stable par $\varinjlim$ filtrantes grandes devant $\pi_0$), et pour tout $\pi \geqslant \pi_0$, $\mathcal{F}^{\pi}$ est stable par les limites filtrantes grandes devant $\pi_0$ et de cardinal $\leqslant \pi_0$ [?].}
\printedpage{58}
% head: A. Grothendieck : Dérivateurs, Ch. XVIII, Édition des Universités de Montpellier II et Paris VII
\item[\textit{c)}] \textit{Si $\pi \geqslant \pi_0$ est} tel que \fbox{$\pi^{\pi_0} = \pi$,} \textit{alors pour tout $F$ dans $\mathcal{M}$, on a un isomorphisme}
\[
F \;=\; \varinjlim_{i \in I} x_i \;,
\]
\textit{les $x_i$ dans $\mathcal{F}^{\pi}$, et $I$ catégorie ordonnée grande devant $\pi_0$. Si de plus $\pi' > \pi$ et $F \in \Ob F^{\pi'}$, on peut prendre $I$ telle que} % [sic] F^{\pi'} (italic F) -> \mathcal{F}^{\pi'}, the filtration is denoted \mathcal{F}^{\pi}
\[
\card I \;\leqslant\; \pi'^{\pi} \qquad \underline{\textit{si}}\ \pi'^{\pi} \geqslant \pi_1. \tag*{(4.7.7.1)}
\]
\item[\textit{d)}] \textit{Si $\pi \geqslant \pi_0$, on a $\mathcal{F}^{\pi} \subset \mathcal{M}_{\pi}$, i.e. les objets de $\mathcal{F}^{\pi}$ sont $\pi$-accessibles.}
\end{enumerate}
\textbf{N.B.} Ce sont là exactement, à part d), les conditions a), b), c) des filtrations cardinales de SGA 4 I 9.7.12, \msp{page 173} à cela près seulement que dans c), on se restreint ici aux ordinaux $\pi$ tels que $\pi^{\pi_0} = \pi$ (\footnote[72]{et \textit{de plus}, dans la majoration (4.7.7.1), aux cardinaux $\pi' > \pi$ tels que l’on ait de plus $\pi'^{\pi} \geqslant \pi_1$ (ce qui est automatique si on a $\pi \geqslant \pi_1$, ce qui est le cas, vu que $\pi = \pi^{\pi_0} \geqslant 2^{\pi_0}$, si on a $2^{\pi_0} \geqslant \pi_1$.)}). La condition supplémentaire d) est une “condition d’accessibilité”, cf. plus bas.

\textsc{Démonstration de} 4.7.7. a) Comme on a
\[
\Filt'^{\pi} \;\subset\; (C^{\wedge})_{\pi} \cap \mathcal{M} \;\subset\; \mathcal{M}_{\pi} \;,
\]
on est réduit à prouver que $\mathcal{M}_{\pi}$ est équivalente à une petite catégorie. On est ramené au cas $\pi \geqslant \pi_1$, $\pi = \pi^{\pi_0}$, auquel cas on a $\mathcal{M}_{\pi} = \Filt^{\pi}$. En fait, on trouve :

\textbf{Corollaire 4.7.7.1} \texttt{[deux fois 4.7.7.1]} (\footnote[73]{Comparer avec 4.7.19.}). \textit{Si $\pi \geqslant \pi_1$, on a}
\[
\left\{
\begin{array}{rcll}
\cardred \Filt^{\pi} & \leqslant & 2^{\pi} & \\[4pt]
\cardred \Filt'^{\pi} & \leqslant & \cardred \mathcal{M}_{\pi} \;\leqslant\; 2^{(\pi^{\pi_0})} & \qquad \textit{($= 2^{\pi}$ si $\pi^{\pi_0} = \pi$).}
\end{array}
\right.
\]

La deuxième majoration résulte de la première et de
\[
\mathcal{F}^{\pi} \;\subset\; \mathcal{M}_{\pi} \;\subset\; \Filt^{\pi^{\pi_0}} \qquad (\text{pour } \pi^{\pi_0} \geqslant \pi_1) \qquad \text{(\footnotemark[74]).}
\]
\footnotetext[74]{Cf. (4.7.6.2).}
Pour la première, on utilise la condition a$''$) de 4.7.1 (qui caractérise $\Filt^{\pi}$ pour $\pi \geqslant \pi_1$). Je passe sur les détails de la vérification.

\msp{page 174}
Prouvons 4.7.2 b). La condition $\mathrm{L}_{\pi_0}$ est connue, reste la stabilité de $\mathcal{F}^{\pi}$. Dans le cas de $\Filt'^{\pi}$, elle est acquise par définition. Dans les cas $C^{\wedge}_{\pi} \cap \mathcal{M}$ et $\mathcal{M}_{\pi}$, elle résulte de la stabilité de la notion de $\pi$-accessibilité par $\varinjlim$ de cardinal $\leqslant \pi$, soit dans $C^{\wedge}$, soit dans $\mathcal{M}$ (compte tenu que $C^{\wedge}$ et $\mathcal{M}$ ont une famille génératrice $C$ formée d’objets $\pi_0$-accessibles, donc $\pi$-accessibles).

Prouvons c) : c’est essentiellement le lemme 4.7.3, compte tenu que l’on a
\[
\Filt^{\pi} \;\subset\; \mathcal{F}^{\pi} \;,
\]

\end{document}
